%!TEX root = ../main.tex
\section{Conclusion}
This research evaluates network encryption and traffic analysis from many points of view; we have demonstrated that even though modern encryption protocols, like TLS and IPsec, are fundamental to data confidentiality, they are not sufficient to guarantee complete protection. 
\\The main vulnerability is what encryption leaves behind and that would be the metadata: it can provide a rich vector of information such as packet sizes, directional flows and timing that an attacker can use to reconstruct user behavior, classify encrypted traffic from applications and decode sensitive communication protocols.
\\\\With the examples examined, we can map the widespread of the vulnerabilities from everyday life with the Android's FBE and Google's GBoard NWP to national defense frameworks with the military encryption discussion. The leak of sensitive information from the metadata affects every data that is on the network, and especially in military networks it is a vulnerability that needs to be assessed. 
Encryption alone is not enough to protect military networks and the leak of operational data can generate dangerous consequences.
\\\\We have come to the conclusion that in order to combat these highly adaptive, AI driven attacks, cyber defense must shift from static models to more dynamic ones. 
As discussed, transforming defense into an automated, proactive mechanism that dynamically adapts to the real-time threat can mitigate the vulnerability.
Lastly its necessary to ensure an integrated, multi-layer approach that not only secures the payload but also masks, obfuscates and dynamically defends the metadata.